\section{Evoluzione del Web e Tecnologie Client-Side}
L'evoluzione del Web client-side si fonda sulla sinergia di tre tecnologie principali: HTML per la struttura, CSS per la presentazione e JavaScript per il comportamento. 

\subsection{Il linguaggio di marcatura HTML}
Il \textit{Hypertext Markup Language} (HTML) è il componente fondamentale del Web; definisce il linguaggio e la struttura dei contenuti disponibili in rete. Lo scopo dell'HTML è quello di strutturare e definire semanticamente le informazioni. La sua evoluzione è contraddistinta da differenti fasi che corrispondono a versioni discrete dello stesso (come l'HTML4 o l'XHTML), ma ad oggi le specifiche sono gestite come uno ``Living Standard'' mantenuto dal gruppo WHATWG (\textit{Web Hypertext Application Technology Working Group}). La particolarità di questo modello prevede un aggiornamento continuo delle specifiche senza l'utilizzo di rilasci numerati \cite{HTMLStandard}.

Alla base del linguaggio c’è il concetto di \textit{Hypertext} (ipertesto), che definisce una struttura di collegamenti (\textit{link}) permettendo la navigazione tra diverse pagine web o all'interno dello stesso documento. L'architettura del \textit{World Wide Web} (WWW) si fonda proprio sulla pubblicazione di contenuti e sulla loro interconnessione tramite questi nodi essenziali della rete.

L'HTML non è un linguaggio di programmazione, bensì un linguaggio di marcatura. Utilizza una sintassi basata su \textit{tag} (etichette racchiuse tra parentesi angolari \texttt{< >}) per istruire il browser su come interpretare e visualizzare i dati. Ad esempio, il \textit{tag} \texttt{<h1>} definisce un'intestazione principale, il \textit{tag} \texttt{<p>} delimita un paragrafo, e il \textit{tag} \texttt{<a>} (denominato \textit{Anchor}) abilita la funzione ipertestuale, permettendo di collegare risorse informative differenti all'interno del WWW.

Quando un browser riceve un documento HTML, ne genera una rappresentazione interna denominata \textit{Document Object Model} (DOM). Il DOM è una struttura gerarchica ad albero che mappa tutti gli elementi della pagina, rendendoli accessibili e manipolabili tramite linguaggi di scripting. Ogni documento segue una struttura ben definita in cui gli elementi possono essere annidati. I contenuti visibili all'utente sono racchiusi all'interno del \textit{tag} \texttt{<body>}, mentre le informazioni di servizio, i metadati o i riferimenti a file esterni risiedono nel \textit{tag} \texttt{<head>}. Infine, per differenziare e raggruppare elementi specifici, l'HTML fornisce attributi globali come \texttt{id} (per identificare in modo univoco un singolo elemento) e \texttt{class} (per classificare gruppi concettuali di elementi), che fungono da fondamentali punti di aggancio per stili e script.

Gli elementi appena descritti permettono all'HTML di definire lo scheletro informativo della pagina, delegando però gli aspetti visivi e comportamentali ad altri linguaggi specializzati \autocite[cap.~5]{niederstrobbinsLearningWebDesign2025}.

\newpage

\subsection{La formattazione tramite CSS} 
\begin{figure}[t]
	\centering
	\begin{lstlisting}[
		language=HTML,
		basicstyle=\small\ttfamily,
		keywordstyle=\color{blue},
		commentstyle=\color{gray},
		stringstyle=\color{purple},
		frame=single,
		breaklines=true
		]
		/* 1. Anatomia di una regola CSS singola */
		h1 {
			color: #2c3e50;
			font-size: 24px;
		}
		
		/* 2. Raggruppamento di piu selettori per ottimizzazione del codice */
		h1, h2, p {
			font-family: 'Helvetica', sans-serif;
			margin-bottom: 15px;
		}
	\end{lstlisting}
	\caption{Esempio di regole CSS: anatomia di una regola e raggruppamento dei selettori.}
	\label{fig:css_regole}
\end{figure}
I \textit{Cascading Style Sheets} (CSS) rappresentano il linguaggio dedicato esclusivamente alla formattazione ed all'aspetto visivo dei documenti \textit{web}. Se il HTML fornisce la semantica e i contenuti, il CSS interviene per garantire una rigorosa ``separazione'' tra i contenuti stessi e la loro presentazione esteriore. \\

Questo paradigma consente il controllo centralizzato dell'aspetto visivo, permettendo di definire in modo totalmente indipendente e granulare dettagli cruciali come i colori, la tipografia, gli spazi (margini e \textit{padding}) e la disposition degli elementi sullo schermo.\\

Il linguaggio CSS si organizza attraverso un sistema di regole che collegano gli elementi di una pagina \textit{web} al loro aspetto estetico. Ogni regola inizia con un selettore, ovvero il termine che indica al \textit{browser} quali parti del testo o della pagina devono essere modificate. Dopo il selettore si trova la dichiarazione, scritta tra parentesi graffe, che contiene i dettagli pratici dello stile. Questa si divide a sua volta in una proprietà, che sceglie la caratteristica da cambiare come il colore o la grandezza, e in un valore, che stabilisce la modifica effettiva. Per rendere il lavoro più veloce e il codice più pulito, è possibile unire più selettori insieme in modo da assegnare lo stesso stile a diversi elementi contemporaneamente con un'unica istruzione.\autocite{CascadingStyleSheets} \\

Uno dei vantaggi più significativi di questa separazione è la possibilità di raccogliere tutte le regole all'interno di file esterni dotati di estensione \texttt{.css}. Un documento HTML può collegarsi a questi file esterni utilizzando il \textit{tag} \texttt{<link>} inserito nel suo \texttt{<head>}. Questa tecnica consente a decine o centinaia di pagine \textit{web} di condividere e puntare al medesimo foglio di stile; di conseguenza, una singola modifica apportata al file CSS si ripercuote istantaneamente sull'intero sito \textit{web}, facilitando scalabilità, coerenza visiva ed operazioni di manutenzione. \autocite[ch. 11]{niederstrobbinsLearningWebDesign2025}\\


\subsection{Framework per il Design Responsivo: Bootstrap}
L'ambiente in cui il \textit{Web} viene fruito oggi è estremamente frammentato: i contenuti devono essere visualizzati su schermi di dimensioni diverse, dai grandi monitor \textit{desktop} ai \textit{tablet}, fino ai piccoli \textit{display} degli \textit{smartphone}, nonché su dispositivi assistivi o degli schermi televisivi. La necessità di risolvere questo problema ha imposto l'adozione del ``\textit{Design} Responsivo'' e di filosofie progettuali come il ``\textit{Mobile First}'', introdotto da specialisti del settore come Luke Wroblewski, che spingono per ottimizzare i siti \textit{web} primariamente per i dispositivi mobili prima di adattarli agli schermi più ampi.\\

L’implementazione di questi paradigmi invece è molto complessa. Sviluppare tali tecnologie da zero, richiede molta precisione ed esperienza da parte dello sviluppatore. Per velocizzare e standardizzare questo processo, vengono ampiamente utilizzati i \textit{framework}, ovvero insiemi organizzati di librerie, convenzioni e codici. \textit{Bootstrap} \autocite{GetStartedBootstrap} integra nativamente sistemi a griglia e regole flessibili, che traducono operativamente i principi del \textit{design} responsivo, riorganizzando dinamicamente i blocchi di contenuto in base alla larghezza e all'orientamento dello schermo (\textit{Viewport}) del dispositivo. \\

Oltre alla gestione del \textit{display}, \textit{Bootstrap} risolve il problema della frammentazione nell’interpretazione degli standard del codice da parte dei vari \textit{browser}, dato che ogni \textit{browser} interpreta gli standard del codice in modo differente. Fogli di stile come \textit{Normalize.css} (integrati in \textit{Bootstrap}) standardizzano i comportamenti di base di elementi grafici, \textit{font} e spaziature. Inoltre integra direttamente i ruoli ARIA, che permettono di garantire la compatibilità con lettori di schermo e ausili per disabilità.\\

In questo panorama, \textit{Bootstrap} si posiziona come uno dei \textit{framework} \textit{frontend} più completi e diffusi al mondo, concepito per consentire una maggiore rapidità di sviluppo.\\


\subsection{Interattività tramite JavaScript} 
Se il HTML è la struttura e il CSS è la presentazione, \textit{JavaScript} è il motore comportamentale del \textit{Web}. \textit{JavaScript} è un linguaggio di programmazione ad alto livello, multi-paradigma (supportando stili imperativi, funzionali e orientati agli oggetti) a tipizzazione dinamica, la cui esecuzione nei moderni motori di \textit{rendering} avviene tipicamente tramite compilazione \textit{Just-In-Time} (JIT).\\

Sebbene nato come linguaggio di \textit{scripting} elementare, si è evoluto nello standard \textit{ECMAScript}. In particolare, il rilascio di ES6 (2015) ha colmato il divario con i linguaggi di livello \textit{enterprise}, introducendo il supporto nativo alla programmazione modulare e una sintassi basata su classi per l'ereditarietà, agevolando l'approccio orientato agli oggetti.\autocite{ECMA26216thEdition}\\

I programmi \textit{JavaScript} necessitano di un ``ambiente \textit{host}'' (\textit{host environment}) per interagire con l’esterno, che nel dominio dello sviluppo \textit{frontend}, è il \textit{browser} \textit{web} stesso. Il \textit{browser} dota \textit{JavaScript} di svariate API (\textit{Application Programming Interface}, interfacce di programmazione esposte all'applicazione). La più importante è il DOM, attraverso la quale il linguaggio può ricercare elementi, mutare l'albero dei nodi, inserire o rimuovere \textit{tag} e modificarne i contenuti.\\

\textit{JavaScript} opera seguendo un'architettura guidata dagli eventi e di tipo asincrono. L’applicazione non impegna la CPU in cicli di attesa attiva, ma risponde a stimoli esterni, definiti eventi. Gli eventi possono essere azioni compiute dall'utente, come un clic o la pressione di un tasto, oppure risposte inviate da un \textit{server}. Per rispondere a questi stimoli, gli sviluppatori utilizzano le funzioni di \textit{callback}, ovvero blocchi di istruzioni che vengono attivati dal \textit{browser} solo nel momento in cui si verifica l'evento a cui sono stati associati. Questo meccanismo, orchestrato dall'interazione tra la coda dei messaggi (\textit{callback queue}) e l’\textit{Event Loop}, garantisce la natura non bloccante del linguaggio. Essendo \textit{JavaScript} un \textit{runtime} \textit{single-threaded}, delegare le operazioni asincrone permette di non occupare il processo (\textit{thread}) principale, preservando la fluidità del \textit{rendering} ed evitando colli di bottiglia nell'esperienza utente.\autocite{flanaganJavaScriptDefinitiveGuide2020}\\

\subsection*{Modelli di Comunicazione Asincrona}
La comunicazione asincrona, storicamente definita con l'acronimo AJAX (\textit{Asynchronous JavaScript and XML}), veniva gestita inizialmente tramite \textit{callback}. Oggigiorno per il fenomeno del ``\textit{callback hell}'' (ovvero l'eccessivo annidamento delle funzioni di ritorno), sono state introdotte le \textit{Promise}, oggetti che rappresentano il risultato futuro e non ancora disponibile di un'operazione. Esse permettono di concatenare in modo lineare le richieste tramite \textit{fetch} e consentono poi di gestire gli errori in modo centralizzato tramite \texttt{catch()}, indirizzando i successi attraverso \texttt{then}.  

Per gestire le \textit{Promise} sono stati introdotti le parole chiave \texttt{async} e \texttt{await}, che permettono di scrivere il codice asincrono, come quello sincrono, sfruttando il medesimo modello non bloccante basato sull'\textit{Event Loop}. \cite{\autocite[ch. 13]{flanaganJavaScriptDefinitiveGuide2020}}

\subsection*{La Libreria jQuery}
\textit{jQuery}\autocite{openjsf.orgJQueryAPIDocumentation} è una libreria \textit{JavaScript} compatta progettata per ottimizzare la manipolazione del DOM, standardizzare la gestione degli eventi e semplificare l'esecuzione delle richieste asincrone (AJAX).

La gestione delle comunicazioni asincrone è centralizzata all’interno del metodo \texttt{jQuery.ajax} (o dell'alias \texttt{\$.ajax()}), il quale permette di configurare ed effettuare richieste HTTP/HTTPS personalizzate. La libreria dispone inoltre di funzioni per i metodi più comuni, come \texttt{jQuery.get()} e \texttt{jQuery.post()} per le chiamate per richiedere e inviare dati al \textit{server} e \texttt{jQuery.getJSON()} che carica e valuta direttamente dati codificati in formato JSON.

Per la gestione del flusso asincrono, \textit{jQuery} implementa inoltre un proprio modello affine allo standard delle \textit{Promise} attraverso l’oggetto \texttt{jQuery.Deferred()}, che fornisce metodi concatenabili come \texttt{.done()} (per il successo), \texttt{.fail()} (per gli errori) e \texttt{.then()} per registrare le \textit{callback} in modo flessibile. 

L’adozione di questa architettura facilita il trasferimento strutturato e coerente dei messaggi in JSON verso il \textit{backend}, ponendosi come uno strato di transizione ideale prima che la richiesta venga presa in carico dai sistemi di \textit{routing} lato \textit{server}.
